
Benchmark selection for hypothesis validation requires identifying which evaluation frameworks provide appropriate metrics, data characteristics, and task formulations. Manual review of benchmark papers is time-consuming, and no systematic method exists to predict future benchmark suitability from design features using historical validation.

Prior work has organized benchmarks into descriptive taxonomies or measured popularity through citation counts. Papers with Code categorizes benchmarks by task type and domain, requiring researchers to manually review each benchmark's documentation. Citation-based approaches measure usage frequency without extracting design constraints that determine hypothesis-benchmark compatibility. Recommendation systems based on collaborative filtering require prior usage data, limiting applicability to novel hypothesis categories.

This work tests whether benchmark design features create systematic constraints on hypothesis testability that persist temporally. Four design features are extracted from benchmark papers: task type, evaluation metrics, data modality, and dataset size. A standardized extraction protocol is developed and validated via inter-rater agreement (Cohen's kappa). Benchmarks are embedded using SentenceBERT and clustered via k-means to discover coverage families. Citation co-occurrence patterns are measured to test whether methods citing benchmarks from the same family exhibit higher citation overlap than random.

The study is organized around four research questions:

\begin{enumerate}
\item Can design features be extracted objectively at scale? (Tests feature extraction protocol)
\item Do coverage families capture meaningful groupings? (Tests clustering quality)
\item Do families predict citation patterns? (Tests main claim and temporal persistence)
\item Can NLP classify citation contexts reliably? (Tests citation evidence extraction)
\end{enumerate}

Experiments on 20 diverse benchmarks show that all four criteria are met. Feature extraction achieves Cohen's kappa 0.917-1.0, clustering produces intra-family similarity 0.748, and citation overlap within families reaches 78.07\%. The primary organizing principle is data modality, not task type, suggesting that metric compatibility creates stronger coverage constraints than task formulation.
